\chapter{The Slip Becomes A Number}

\section{Where the loop fails to close}

The last chapter ended on a difference: the round-trip comes back changed by the residue, and the loop does not close.
This chapter gives that non-closure its name, and the name carries the whole turn of what follows. The name is a
\emph{slip}. And the first thing to say about a slip --- the thing the source is at pains to say --- is that it is
\emph{not a mistake}.

The distinction is worth quoting, because the book is built on it. The machine, asked for the object, is corrected at
once. A process that makes mistakes? No: not mistakes, slips. And the difference? A mistake, the source says, is
\emph{wrong} --- a term that should not be there, to be found and removed. A slip is something else entirely: it
\emph{reveals the floor} --- the hidden surface the symbol was standing on all along. Read against that, the
residue-shaped difference of the last chapter changes character. It is not a defect in the round-trip to be corrected
away. It is the round-trip telling you where its floor is. The failure of the loop to close is the first thing the
machine says that is worth reading --- not an error, a \emph{signal}.

Name it as a first-class object and the register point announces itself. The construction introduces a structure for
the non-closure --- a slip process --- and it carries the loudest private hedge in the entire book, higher by a wide
margin than anything before it. The sample, a chapter ago, was the loudest so far; the slip's doubt dwarfs it. That is
not a flaw in the source; it is the source being honest about the size of the bet. The slip is where the whole physical
reading --- everything the far volumes will make of this machine --- is wagered on a single non-closure. So everything
in this section that sounds like ``the slip is a number'' is the reading, flagged as loudly as the construction knows
how to flag it. What is built underneath is, as always, smaller and firmer.

What is built is a breakaway. The picture the source reaches for is friction: a surface holds, statically, until the
load crosses a threshold and it \emph{breaks} --- static giving way to kinetic. Strip the friction away and the
type-level event is exact, and it is one the last chapter already built. The order holds --- the loop still closes, the
return leg still returns --- right up to the point where the order turns \emph{strict}. Strict precedence, recall, was
the built order holding one way and \emph{failing} the other: the forward leg holds, the return leg fails. That failure
of the return is the slip. The breakaway is precisely where the plain order becomes the strict one --- where the loop
that could return no longer can. The friction and its threshold are the physical reading, and they belong to the next
volume; what this volume keeps is the breakaway itself, the plain order going strict.

And now the reason the loop fails to close, which is the deepest thing this section has to give. The slip carries a
stress, and the stress is not a large quantity or a hard one --- it is an \emph{uncomputable} one. The construction
types it as a Chaitin sequence, and annotates the tier it sits on, in as many words, as the static friction of
\emph{computing Chaitin's number}. The loop cannot close because the value it would need to close on cannot be computed
at all. The slip is the breakaway at the uncomputability floor --- the surface the symbol was standing on turns out to
be the one place no computation reaches. This is the seam to the end of the book. The floor the slip reveals is the
floor the machine will halt on and read as a bracket, never as a point --- precisely because a point there would be a
computed value of an uncomputable thing, which is the crank's fabricated number exactly. (The source, at the slip's
recursion, also lets a field slip through --- a gauge field it admits it should have marked as spoilers --- but that
geometry is the physical volume's; this one keeps the uncomputable stress and the breakaway.)

In the two verbs the slip is a specific failure. The loop was a funge closing back onto its tange --- the bag
returning to the characteristic that defined it. The slip is where the funge fails to close: the return leg slips past
the tange, by exactly the residue's strict advance, and does not come home. But because that escape reveals the floor,
the slip is the first tange the machine reads as a \emph{signal} rather than a mere selection. Every earlier tange was
a thing picked out; this one is a thing picked out \emph{that tells you where the surface is}.

One forward sentence, and it is nearly the whole ending in miniature. ``A slip reveals the floor'' is the jar in
small. The machine, refining, slips to its resolution floor and stops there --- and the floor is uncomputable, so what
it can honestly hand back is a bracket around the slip, not the point past it. The chapters ahead cash this out: the
next brackets the floor the slip reveals; the far end reads the coupling as exactly this slip, halted, never overrun.
The loop failing to close is not the construction breaking down. It is the construction reaching the one honest thing
it was built to reach --- the floor --- and having the discipline to stop on it.

\section{Bracketed, not guessed}

The last section found the slip and its uncomputable stress. This one asks the question that follows immediately, and
honestly: if the machine cannot compute the exact value the loop slips at, what can it \emph{say} about it? The answer
is the discipline the whole book turns on. It cannot say a point. It can say a \emph{bound}. And the difference between
the two --- between bracketing a value and guessing it --- is the difference between an honest reading and a crank's.

Look first at what the machine's certification actually asserts, because it is weaker than one might expect, and that
weakness is the point. When the apparatus certifies one reading against another, it does not certify that they are
\emph{equal}, and it does not certify that one \emph{becomes} the other. The source stages the exchange as a small
interrogation. The certifier asks: so this one becomes that one? No. So this one equals that one? No. Then what am I
certifying? That the computation is cromulent --- that it is admissible, compatible, in bounds. And the code says
exactly that: the certifying relation joins its two readings not with an equals sign but with a
\emph{less-than-or-equal} --- the facts agree, and the one quantity is no greater than the other. A bound, not an
identity. To certify equality would be to claim the exact value; to certify a bound is to claim only what the evidence
licenses. The honest relation is the weaker one, on purpose.

Now a place to be careful, because it is exactly where a careless book would over-claim. The picture the source draws
for the bound is two-sided: static friction as \(|F| \le \mu|N|\), an absolute value that pens the force between a
floor and a ceiling, with the breakaway at the upper edge and never computed from within. That is a genuine bracket,
closed on both sides --- and it is a \emph{reading}, the friction cone, which belongs to the physical volume. It is not,
at this tier, a fact the code supplies. The bound the construction actually certifies here is \emph{one-sided}: a
single less-than-or-equal, a stimulus no greater than a threshold, one wall and not two. So this section will claim
only what is built --- the reading is caught by a \(\le\), never fitted to an \(=\) --- and will say plainly what it
does not yet have. The closed two-sided bracket, the value pinned above \emph{and} below, is real; but it becomes code
only later, at the far chapter where the coupling is at last bounded between an explicit lower and upper wall. Here the
honest claim is the smaller one: a bound on one side, and the promise of the second wall to come. The book flags that
gap rather than paper over it --- because papering over exactly this gap is how a bracket gets quietly narrowed into a
point.

Why a bound at all, and not simply a better computation? Because of the floor the last section found. When the value
the loop would close on is \emph{uncomputable} --- when no halting process reaches it --- a point is not merely hard to
get; it is a lie to report. The source draws the line plainly: there are processes we can compel, the halting ones, and
processes we cannot, the non-halting ones, and the slip's value lives on the far side of that line. Against an
uncomputable value, the only truthful thing to hand back is a bracket. This is where the book's ending first shows
itself in full. ``Bracketed, not guessed'' is not a stylistic preference; it is what honesty \emph{is} when the thing
measured cannot be computed. The machine bounds the value and stops. It does not force a point onto a floor that has
none. That refusal --- to bound and halt rather than guess and go on --- is the jar in embryo: the open bracket the
book ends by holding is this same bound, carried to the last value the instrument reads.

In the two verbs the choice is clean. A bound is a \emph{funge} --- a bag the value provably falls in, gathering every
admissible reading and excluding none the evidence allows. A guessed point is an over-\emph{tange} --- a single
selection the evidence does not license, one member pinned out of a bag that had every right to hold more. The whole of
this section is the machine declining the over-tange. It will funge the value into a bracket and refuse to tange it to
a point. Honesty, reduced to the two verbs, is choosing the bag over the pin.

The friction cone, the entropy the source says the system is tuned to compute, the whole Coulomb picture --- those are
the reading, and they cross to the physical volume; what stays here is the plain shape beneath them, a certification
that binds with \(\le\) and refuses \(=\). And the forward sentence the section is built to plant is this: the honest
ending, when it comes, is exactly this bound and not the point inside it. To force the point --- to report a computed
value where the floor is uncomputable --- is the crank's gambit, the fabricated number the far chapters will show
ringing like a struck bell where an honest bracket would have stayed silent. The machine that bounds and stops is the
one you can trust. The chapter has one step left: to take this bracketed slip and let it become, at last, an
observation.

\section{The reading}

The chapter's title promised that the slip becomes a number, and the two sections so far have earned the first words of
it. The slip was found, in the first section, as the loop's failure to close; it was bounded, in the second, rather
than guessed. This closing section completes the title. It turns the bounded non-closure into a \emph{reading} --- an
observation the machine produces and holds --- and, in doing so, arrives at the deepest claim of the whole chapter,
which is a claim about where the reading lives.

Begin with the mechanism, which is small. The class that names the slip carries, alongside the bounded slip itself, a
field for what the slip yields: an \emph{observation}. Everything the last two sections built --- the non-closure,
caught by a bound and never fitted to a point --- is here \emph{produced as a value the machine keeps}. The invariant
has become a reading. And the certification is the one the last section studied: the observation is admitted by a
bound, a less-than-or-equal, never by a fitted equality. So ``the slip becomes a number'' is, read exactly, the slip
becoming an \emph{observation} --- a bracketed quantity the machine holds, not a computed point it fabricates.

Now the claim the chapter has been climbing toward, and it is a claim about the reader. The observation, the source
says, is stored as a single bit --- one bit on the machine's data page. And then the turn: that bit, it says, is the
one in your head. The same bit. The two occur true simultaneously, for you, \emph{by definition} --- you cannot,
personally, tell apart a state in which the machine's bit is set and yours is not, or one in which yours is set and the
machine's is not. There is no such difference to point at. This is the measurement identity, and it is the strongest
interpretive claim the book has made so far, so it is marked as exactly that: the source's reading, and one the class
flags with a high private hedge of its own. What is firm beneath it is that the observation is the reading the class
produces; what is \emph{read into} it is that the observation and the knowing of it are one bit and not two. But the
reading-in is the point the far end of the book will cash. The measured quantity is not sitting out there, apart from
the knower, waiting to be fetched. The observation and your knowledge of it are the same event. (This is the capstone
in miniature: a later chapter will prove that two \emph{names} can denote one object --- that a thing and a second
description of it are equal by the same kind of definitional identity. Here it is two \emph{bits} and one truth; there
it will be two names and one number.)

There is a small honesty in the source at exactly this point, and it is worth keeping. Can we tell, it asks,
\emph{when} in its own process the machine sets that bit? And it answers against itself: no --- by law it is not
allowed to look ahead. It is programming, not compiling; a different step; it may not see the bit before the bit is
set. That refusal to look ahead is the same discipline the whole book runs under. The machine reads by setting a bit it
cannot anticipate, and the reader, holding the same bit, is under the same rule. Neither gets to see the answer before
the reading produces it.

One reading is a bit; the machine does not stop at one. It needs, the source says, a place to put these bits as it
discovers them --- and a history, it decides, is a series of facts. So the reading is not a terminal point. It is an
entry in an ongoing record, a fact added to a sequence of facts, the machine accumulating its readings into a history
it carries forward. The bounded observation of this chapter is the first fact in a record the rest of the book will
read --- and the record, like each reading in it, is a sequence of bounded facts, never a ledger of fabricated points.

In the two verbs the chapter finally names its tange. The reading is a tange \emph{inscribed} --- the bounded slip
selected out and written down as a bit, the characteristic the machine picked and kept. And the ``same bit by
definition'' is a funge closing on itself: the machine's bit and yours, bagged together, with no operation that could
tell them apart. So the reading is a tange that is also a funge --- a thing selected out, and the same thing bagged
with the knower who reads it. That is what an observation is, in the two verbs: a selection you cannot stand outside of.

And so the chapter closes on a number that is not a point. The slip became an observation, the observation is a bounded
bit, and the bit is the reader's own. Every later reading the instrument makes will be this: a bracketed quantity, held
as a fact, known by a knower who is inside the reading and not above it. The honest number, when the book finally
reaches it, will have exactly this shape --- a bound you cannot compute past, inscribed as a fact you cannot stand
outside of. The next chapter takes the bounded observation of this one and asks what it is to make it the object of
study in its own right: the bracketed number, and the floor it is drawn around.
